% 第15批：《大学物理学学习指导》自测题2选择2.1.1—2.1.9。
% 印刷49—51页／PDF物理57—59页；逐题按扫描页重述、推导。
\subsection{自测题2：质点动力学和刚体力矩选择}

\sourcetest{2.1.1}{49}{57}
\problemtext 质点在平面内的位矢$\bm r=at^2\bm i+bt^2\bm j$（$a,b$常量），判别路径与运动性质。
\workings
两个坐标分量成固定比例$y/x=b/a$（$a\ne0$），轨迹为直线；
$\bm v=2t(a\bm i+b\bm j)$随时间改变，因此是\textbf{变速直线运动}，选\textbf{B}。
\crosscheck $a=0$时轨迹为$y$轴，同样是直线；不能将位矢有两个分量误判为曲线运动。

\sourcetest{2.1.2}{49}{57}
\problemtext 从加速度方向、平均速率与平均速度、恒定速率和速度变化四条说法中选出正确者。
\workings
恒定加速度不保证运动方向恒定，例如平抛运动；
平均速率是路程除时间，平均速度是位移除时间，两者并不恒相等；
一般变速运动的平均速度不能直接由始末速度算术平均求得。
恒定\emph{速率}仍可改变速度方向，例子为匀速圆周运动，
因此选\textbf{D}。
\crosscheck 区分标量速率与矢量速度。

\sourcetest{2.1.3}{49}{57}
\problemtext 与水平成30°的斜向上力$F$把重为$G$的物块压靠竖直墙。
若无论增大$F$多少均不能使物块向上滑动，求静摩擦系数$\mu$的条件。
\workings
法向力$N=F\cos30^\circ$，
切向向上的力为$F\sin30^\circ$。
物块不向上滑的必要条件
$F\sin30^\circ\le G+\mu F\cos30^\circ$；
令$F$增大而$G$固定，得
\[
\boxed{\mu\ge\tan30^\circ=\frac1{\sqrt3}}.
\]
选\textbf{B}。
\crosscheck 不能把重力视作随$F$增大而同步增加的摩擦能力。

\sourcetest{2.1.4}{49}{57}
\problemtext 质量$m_B=2m_A$的物块A、B以压缩弹簧连接，
在光滑水平桌上从静止释放后，求动能$E_{kA}/E_{kB}$。
\workings
系统初始总动量0，释放后$p_A=-p_B$，
因此
\[
\boxed{\frac{E_{kA}}{E_{kB}}
=\frac{p^2/(2m_A)}{p^2/(2m_B)}
=\frac{m_B}{m_A}=2}.
\]
选\textbf{B}。
\crosscheck 动量大小相同、质量较小的一方动能较大。

\sourcetest{2.1.5}{49}{57}
\problemtext 甲乙体重身高相同，各攀同一无摩擦定滑轮轻绳一端，
均从静止开始；虽然甲相对绳的爬速是乙的两倍，问谁先到顶。
\workings
两端轻绳具有相同张力$T$。
若二人质量均为$m$，对地加速度都为$a=T/m-g$；
同初始位置与初速则对地速度与位移始终相同，
\[
\boxed{y_A(t)=y_B(t)}.
\]
因此\textbf{同时到达}，选\textbf{C}。
\crosscheck “相对绳子爬得快”不代表“对地上升快”，
轻绳在两侧的运动方向相反。

\sourcetest{2.1.6}{49--50}{57--58}
\problemtext 源图2.1.3：斜面弹簧劲度系数$k$，
物块初始位置$a$较自然长度端点$O$靠左$x_1$，
沿坡移动至$b$的位移为$x_2$；求弹簧、物块和地球系统的势能变化。
\workings
初态弹簧压缩量$x_1$；末态弹簧有向形变量$x_2-x_1$。
物块上升$x_2\sin\alpha$，故
\[
\boxed{\Delta U
=\tfrac12 k(x_2-x_1)^2-\tfrac12kx_1^2
+mgx_2\sin\alpha}.
\]
选\textbf{C}。
\crosscheck 计算变化量必须减去原有弹性势能，不可只写末态势能。

\sourcetest{2.1.7}{50}{58}
\problemtext 有转动惯量的定滑轮两端挂$m_1<m_2$，
虽此刻轮子逆时针转动，比较左、右绳张力。
\workings
较重的$m_2$使系统角加速度为顺时针，尽管瞬时角速度可逆时针。
设左张力$T_1$、右张力$T_2$，对轮轴
$(T_2-T_1)R=I\alpha_{\rm clockwise}>0$，故
\[
\boxed{T_2>T_1}.
\]
选\textbf{C（右边大）}。
\crosscheck 角速度方向与角加速度方向不必相同。

\sourcetest{2.1.8}{50}{58}
\problemtext 两个力作用于绕定轴的刚体，
判断“均平行转轴”“均垂直转轴”“合力为零”“对轴合力矩为零”
四种条件与轴向力矩之间的关系。
\workings
均平行转轴时各自轴向力矩均为0，命题(1)真。
均垂直转轴时两轴向力矩\emph{可能}抵消，命题(2)真。
合力为零的\emph{力偶}仍可有力矩，命题(3)假。
合力矩为零不要求合力为零，命题(4)假。
故选\textbf{B（1、2正确）}。
\crosscheck 题问的是对\emph{固定轴}的力矩，不是完整三维力矩向量。

\sourcetest{2.1.9}{50--51}{58--59}
\problemtext 质量$M$、长度$L$的均匀细棒以左端O为转轴，
$J_O=ML^2/3$。质量$m$子弹以垂直于杆的$v$
穿过自由端，穿出后仍同向以$v/2$飞行，
求碰后杆角速度。
\workings
短暂穿透期间系统关于O角动量守恒，子弹损失
$mLv-mL(v/2)=mLv/2$，
\[
\frac13ML^2\omega=\frac12mvL,
\qquad
\boxed{\omega=\frac{3mv}{2ML}}.
\]
选\textbf{B}。
\crosscheck 子弹穿透一般耗散机械能，不能使用机械能守恒替代角动量守恒。
